\section{Determinants: Basic Properties}

\subsection{Conventions}

\begin{convention}
\label{conv.det.K}
For the rest of this section, we fix a commutative ring $K$.
In most examples, $K$ will be $\mathbb{Z}$ or $\mathbb{Q}$ or a polynomial ring.
\end{convention}

\begin{convention}
\label{conv.det.matrices}
Let $n, m \in \mathbb{N}$.

\textbf{(a)} If $A$ is an $n \times m$-matrix, then $A_{i,j}$ shall mean the
$(i,j)$-th entry of $A$, that is, the entry of $A$ in row $i$ and column $j$.

\textbf{(b)} If $a_{i,j}$ is an element of $K$ for each $i \in [n]$ and each
$j \in [m]$, then
\[
  \left( a_{i,j} \right)_{1 \leq i \leq n,\; 1 \leq j \leq m}
\]
shall denote the $n \times m$-matrix whose $(i,j)$-th entry is $a_{i,j}$ for
all $i \in [n]$ and $j \in [m]$.

\textbf{(c)} We let $K^{n \times m}$ denote the set of all $n \times m$-matrices
with entries in $K$. This is a $K$-module. If $n = m$, this is also a $K$-algebra.

\textbf{(d)} Let $A \in K^{n \times m}$ be an $n \times m$-matrix. The
\emph{transpose} $A^T$ of $A$ is defined to be the $m \times n$-matrix whose
entries are given by
\[
  \left( A^T \right)_{i,j} = A_{j,i}
  \quad \text{for all } i \in [m] \text{ and } j \in [n].
\]
\end{convention}

\subsection{Definition}

\begin{definition}
\label{def.det.det}
\lean{AlgebraicCombinatorics.Det.det_eq_sum_sign_prod_textbook}
\leantarget
\leanok
Let $n \in \mathbb{N}$. Let $A \in K^{n \times n}$ be an $n \times n$-matrix.
The \emph{determinant} $\det A$ of $A$ is defined to be the element
\[
  \sum_{\sigma \in S_n} (-1)^{\sigma}
  \underbrace{A_{1,\sigma(1)} A_{2,\sigma(2)} \cdots A_{n,\sigma(n)}}_{
    = \prod_{i=1}^{n} A_{i,\sigma(i)}}
\]
of $K$. Here:
\begin{itemize}
\item we let $S_n$ denote the $n$-th symmetric group (i.e., the group of
permutations of $[n] = \{1, 2, \ldots, n\}$);
\item we let $(-1)^{\sigma}$ denote the sign of the permutation $\sigma$
(as defined in Definition~\ref{def.perm.sign}).
\end{itemize}
\end{definition}

\subsection{Small determinants and basic API}

\begin{lemma}
\label{lem.det.fin-zero}
\lean{AlgebraicCombinatorics.Det.det_fin_zero'}
\leanhelper
\leanok
The determinant of a $0 \times 0$ matrix is $1$ (empty product convention).
\end{lemma}

\begin{proof}
\leanok
The sum over $S_0 = \{\mathrm{id}\}$ gives $(-1)^{\mathrm{id}} \cdot 1 = 1$
(empty product is $1$).
\end{proof}

\begin{lemma}
\label{lem.det.fin-one}
\lean{AlgebraicCombinatorics.Det.det_fin_one'}
\leanhelper
\leanok
The determinant of a $1 \times 1$ matrix is its single entry: $\det (a) = a$.
\end{lemma}

\begin{proof}
\leanok
The sum over $S_1 = \{\mathrm{id}\}$ gives $(-1)^{\mathrm{id}} \cdot a = a$.
\end{proof}

\begin{lemma}
\label{lem.det.fin-two}
\lean{AlgebraicCombinatorics.Det.det_fin_two'}
\leanhelper
\leanok
The determinant of a $2 \times 2$ matrix satisfies
$\det \begin{pmatrix} a & b \\ c & d \end{pmatrix} = ad - bc$.
\end{lemma}

\begin{proof}
\leanok
$S_2$ has two elements: the identity (sign $+1$) and the transposition (sign $-1$).
The sum gives $a \cdot d - b \cdot c$.
\end{proof}

\begin{lemma}
\label{lem.det.one-eq-one}
\lean{AlgebraicCombinatorics.Det.det_one_eq_one}
\leanhelper
\leanok
The determinant of the identity matrix is $1$: $\det I_n = 1$.
\end{lemma}

\begin{proof}
\leanok
By the fact that the identity matrix has determinant $1$.
\end{proof}

\begin{lemma}
\label{lem.det.zero-eq-zero}
\lean{AlgebraicCombinatorics.Det.det_zero_eq_zero}
\leanhelper
\leanok
For $n \geq 1$, the determinant of the zero matrix is $0$.
(For $n = 0$, the zero matrix is the empty matrix, and $\det(\emptyset) = 1$.)
\end{lemma}

\begin{proof}
\leanok
The zero matrix has a zero row, so $\det = 0$ by
Theorem~\ref{thm.det.rowop}~\textbf{(b)}.
\end{proof}

\begin{lemma}
\label{lem.det.diagonal}
\lean{AlgebraicCombinatorics.Det.det_diagonal'}
\leanhelper
\leanok
The determinant of a diagonal matrix equals the product of its diagonal entries:
$\det(\mathrm{diag}(d_1, \ldots, d_n)) = d_1 d_2 \cdots d_n$.
\end{lemma}

\begin{proof}
\leanok
A diagonal matrix is both upper- and lower-triangular, so this follows from
Theorem~\ref{thm.det.triang}.
\end{proof}

\begin{lemma}
\label{lem.det.const-matrix-zero}
\lean{AlgebraicCombinatorics.Det.det_const_matrix_eq_zero}
\leanhelper
\leanok
For $n \geq 2$, a matrix with all entries equal to $c$ has determinant $0$.
This is a corollary of Proposition~\ref{prop.det.xiyj} with $x_i = 1$ and $y_j = c$.
\end{lemma}

\begin{proof}
\leanok
The constant matrix $(c)_{i,j}$ is the outer product of $\mathbf{1}$ and $(c, c, \ldots, c)$.
By Proposition~\ref{prop.det.xiyj}, its determinant is $0$.
\end{proof}

\begin{lemma}
\label{lem.det.sum-sign-zero}
\lean{AlgebraicCombinatorics.Det.sum_sign_eq_zero}
\leanhelper
\leanok
Let $n \geq 2$. Then,
\[
  \sum_{\sigma \in S_n} (-1)^{\sigma} = 0.
\]
\end{lemma}

\begin{proof}
\leanok
The proof uses a sign-reversing involution: for $n \geq 2$, pick two distinct
elements $i, j$ and their transposition $t_{i,j}$. The map
$\sigma \mapsto \sigma \cdot t_{i,j}$ pairs up permutations with opposite signs,
so all summands cancel.
\end{proof}

\begin{proposition}
\label{prop.det.xiyj}
\lean{AlgebraicCombinatorics.Det.det_outerProduct_eq_zero}
\leantarget
\leanok
Let $n \in \mathbb{N}$ be such that $n \geq 2$. Let
$x_1, x_2, \ldots, x_n \in K$ and $y_1, y_2, \ldots, y_n \in K$. Then,
\[
  \det\left( \left( x_i y_j \right)_{1 \leq i \leq n,\; 1 \leq j \leq n} \right) = 0.
\]
\end{proposition}

\begin{proof}
\leanok
The definition of the determinant yields
(using Lemma~\ref{lem.det.sum-sign-zero} for the last step)
\begin{align*}
  \det\left( \left( x_i y_j \right)_{1 \leq i \leq n,\; 1 \leq j \leq n} \right)
  &= \sum_{\sigma \in S_n} (-1)^{\sigma}
     \underbrace{(x_1 y_{\sigma(1)})(x_2 y_{\sigma(2)}) \cdots (x_n y_{\sigma(n)})}_{
       = (x_1 x_2 \cdots x_n)(y_{\sigma(1)} y_{\sigma(2)} \cdots y_{\sigma(n)})} \\
  &= \sum_{\sigma \in S_n} (-1)^{\sigma}
     (x_1 x_2 \cdots x_n)
     \underbrace{(y_{\sigma(1)} y_{\sigma(2)} \cdots y_{\sigma(n)})}_{\substack{
       = y_1 y_2 \cdots y_n \\
       \text{(since } \sigma \text{ is a bijection)}}} \\
  &= (x_1 x_2 \cdots x_n)(y_1 y_2 \cdots y_n)
     \underbrace{\sum_{\sigma \in S_n} (-1)^{\sigma}}_{= 0} = 0.
\end{align*}
\end{proof}

\begin{lemma}
\label{lem.det.hollow-core}
\lean{AlgebraicCombinatorics.Det.det_hollow_core_eq_zero}
\leanhelper
\leanok
Let $A$ be a $5 \times 5$ matrix such that $A_{i,j} = 0$ whenever $i, j \in \{2, 3, 4\}$.
Then $\det A = 0$.
This formalizes the $5 \times 5$ hollow matrix example: a matrix with a
$3 \times 3$ block of zeros in the middle has determinant $0$.
\end{lemma}

\begin{proof}
\leanok
By the pigeonhole principle: any permutation $\sigma$ of $[5]$
must map some element of $\{2,3,4\}$ to an element of $\{2,3,4\}$, because
$|\{2,3,4\}| = 3 > 2 = |\{1,5\}|$.
Hence, for every $\sigma$, the product $\prod_i A_{i,\sigma(i)}$ contains a zero factor,
making every term in the determinant sum equal to $0$.
\end{proof}

\subsection{Helpers for Proposition~\ref{prop.det.xi+yj}}

\begin{lemma}
\label{lem.det.sumMatrix-factorA-zero-col}
\lean{AlgebraicCombinatorics.Det.sumMatrix_factorA_has_zero_col}
\leanhelper
\leanok
The factor matrix $A$ in the factorization of the sum matrix $(x_i + y_j)$
has a zero column (column index $2$) when $n \geq 3$.
The matrix $A$ has $x_i$ in the first column, $1$ in the second column,
and $0$ elsewhere.
\end{lemma}

\begin{proof}
\leanok
By direct computation: for column index $2$, neither $j = 0$ nor $j = 1$, so the entry is $0$.
\end{proof}

\begin{lemma}
\label{lem.det.det-sumMatrix-factorA-zero}
\lean{AlgebraicCombinatorics.Det.det_sumMatrix_factorA_eq_zero}
\leanhelper
\leanok
For $n \geq 3$, the factor matrix $A$ has determinant $0$.
\end{lemma}

\begin{proof}
\leanok
By Lemma~\ref{lem.det.sumMatrix-factorA-zero-col}, column $2$ is zero.
By Theorem~\ref{thm.det.colop}~\textbf{(b)}, $\det A = 0$.
\end{proof}

\begin{lemma}
\label{lem.det.sumMatrix-eq-factorA-mul-factorB}
\lean{AlgebraicCombinatorics.Det.sumMatrix_eq_factorA_mul_factorB}
\leanhelper
\leanok
For $n \geq 2$, the matrix $(x_i + y_j)_{i,j}$ factors as $A \cdot B$ where:
\[
A_{i,j} = \begin{cases} x_i & \text{if } j = 0, \\ 1 & \text{if } j = 1, \\ 0 & \text{otherwise,} \end{cases}
\qquad
B_{i,j} = \begin{cases} 1 & \text{if } i = 0, \\ y_j & \text{if } i = 1, \\ 0 & \text{otherwise.} \end{cases}
\]
\end{lemma}

\begin{proof}
\leanok
Direct matrix multiplication: $(AB)_{i,j} = \sum_k A_{i,k} B_{k,j} = x_i \cdot 1 + 1 \cdot y_j = x_i + y_j$.
\end{proof}

\begin{proposition}
\label{prop.det.xi+yj}
\lean{AlgebraicCombinatorics.Det.det_sumMatrix_eq_zero}
\leantarget
\leanok
Let $n \in \mathbb{N}$ be such that $n \geq 3$. Let
$x_1, x_2, \ldots, x_n \in K$ and $y_1, y_2, \ldots, y_n \in K$. Then,
\[
  \det\left( \left( x_i + y_j \right)_{1 \leq i \leq n,\; 1 \leq j \leq n} \right) = 0.
\]
\end{proposition}

\begin{proof}
\leanok
Define two $n \times n$-matrices
\[
  A := \begin{pmatrix}
    x_1 & 1 & 0 & \cdots & 0 \\
    x_2 & 1 & 0 & \cdots & 0 \\
    x_3 & 1 & 0 & \cdots & 0 \\
    \vdots & \vdots & \vdots & \ddots & \vdots \\
    x_n & 1 & 0 & \cdots & 0
  \end{pmatrix}
  \quad \text{and} \quad
  B := \begin{pmatrix}
    1 & 1 & 1 & \cdots & 1 \\
    y_1 & y_2 & y_3 & \cdots & y_n \\
    0 & 0 & 0 & \cdots & 0 \\
    \vdots & \vdots & \vdots & \ddots & \vdots \\
    0 & 0 & 0 & \cdots & 0
  \end{pmatrix}.
\]
Then $\left( x_i + y_j \right)_{1 \leq i \leq n,\; 1 \leq j \leq n} = AB$, so
\[
  \det\left( \left( x_i + y_j \right)_{1 \leq i \leq n,\; 1 \leq j \leq n} \right)
  = \det(AB) = \det A \cdot \det B
\]
(by Theorem~\ref{thm.det.detAB}).
However, the matrix $A$ has a zero column
(since $n \geq 3$), and thus $\det A = 0$
(by Theorem~\ref{thm.det.colop}~\textbf{(b)}).
Hence the product is $0$.
\end{proof}

\subsection{Basic properties}

\begin{theorem}[Transposes preserve determinants]
\label{thm.det.transp}
\lean{AlgebraicCombinatorics.Det.det_transpose'}
\leantarget
\leanok
Let $n \in \mathbb{N}$. If $A \in K^{n \times n}$ is any $n \times n$-matrix,
then $\det(A^T) = \det A$.
\end{theorem}

\begin{proof}
\leanok
See \cite[Corollary B.16]{Strick13} or \cite[Exercise 6.4]{detnotes} or
\cite[\S 5.3.2]{Laue-det}.
\end{proof}

\begin{theorem}[Determinants of triangular matrices]
\label{thm.det.triang}
\lean{AlgebraicCombinatorics.Det.det_upperTriangular, AlgebraicCombinatorics.Det.det_lowerTriangular}
\leantarget
\leanok
Let $n \in \mathbb{N}$. Let $A \in K^{n \times n}$ be a triangular (i.e.,
lower-triangular or upper-triangular) $n \times n$-matrix. Then, the
determinant of the matrix $A$ is the product of its diagonal entries. That is,
\[
  \det A = A_{1,1} A_{2,2} \cdots A_{n,n}.
\]
\end{theorem}

\begin{proof}
\leanok
See \cite[Proposition B.11]{Strick13} or \cite[Exercise 6.3 and the paragraph
after Exercise 6.4]{detnotes}.
\end{proof}

\begin{theorem}[Row operation properties]
\label{thm.det.rowop}
\lean{AlgebraicCombinatorics.Det.det_swap_rows, AlgebraicCombinatorics.Det.det_zero_row, AlgebraicCombinatorics.Det.det_eq_rows, AlgebraicCombinatorics.Det.det_scale_row, AlgebraicCombinatorics.Det.det_add_row, AlgebraicCombinatorics.Det.det_add_smul_row, AlgebraicCombinatorics.Det.det_row_add}
\leantarget
\leanok
Let $n \in \mathbb{N}$. Let $A \in K^{n \times n}$ be an $n \times n$-matrix. Then:

\textbf{(a)} If we swap two rows of $A$, then $\det A$ gets multiplied by $-1$.

\textbf{(b)} If $A$ has a zero row (i.e., a row that consists entirely of zeroes),
then $\det A = 0$.

\textbf{(c)} If $A$ has two equal rows, then $\det A = 0$.

\textbf{(d)} Let $\lambda \in K$. If we multiply a row of $A$ by $\lambda$
(that is, we multiply all entries of this one row by $\lambda$, while leaving
all other entries of $A$ unchanged), then $\det A$ gets multiplied by $\lambda$.

\textbf{(e)} If we add a row of $A$ to another row of $A$ (that is, we add
each entry of the former row to the corresponding entry of the latter), then
$\det A$ stays unchanged.

\textbf{(f)} Let $\lambda \in K$. If we add $\lambda$ times a row of $A$ to
another row of $A$ (that is, we add $\lambda$ times each entry of the former
row to the corresponding entry of the latter), then $\det A$ stays unchanged.

\textbf{(g)} Let $B, C \in K^{n \times n}$ be two further $n \times n$-matrices.
Let $k \in [n]$. Assume that
\[
  (\text{the } k\text{-th row of } C)
  = (\text{the } k\text{-th row of } A)
  + (\text{the } k\text{-th row of } B),
\]
whereas each $i \neq k$ satisfies
\[
  (\text{the } i\text{-th row of } C)
  = (\text{the } i\text{-th row of } A)
  = (\text{the } i\text{-th row of } B).
\]
Then,
\[
  \det C = \det A + \det B.
\]
\end{theorem}

\begin{proof}
\leanok
\textbf{(a)} See \cite[Exercise 6.7 \textbf{(a)}]{detnotes}. This is also a
particular case of \cite[Corollary B.19]{Strick13}.

\textbf{(b)} See \cite[Exercise 6.7 \textbf{(c)}]{detnotes}. This is also
near-obvious from Definition~\ref{def.det.det}.

\textbf{(c)} See \cite[Exercise 6.7 \textbf{(e)}]{detnotes} or
\cite[\S 5.3.3, property (iii)]{Laue-det}.

\textbf{(d)} See \cite[Exercise 6.7 \textbf{(g)}]{detnotes} or
\cite[\S 5.3.3, property (ii)]{Laue-det}.

\textbf{(f)} See \cite[Exercise 6.8 \textbf{(a)}]{detnotes}.

\textbf{(e)} This is the particular case of part \textbf{(f)} for $\lambda = 1$.

\textbf{(g)} See \cite[Exercise 6.7 \textbf{(i)}]{detnotes} or
\cite[\S 5.3.3, property (i)]{Laue-det}.
\end{proof}

\begin{theorem}[Column operation properties]
\label{thm.det.colop}
\lean{AlgebraicCombinatorics.Det.det_swap_cols, AlgebraicCombinatorics.Det.det_zero_col, AlgebraicCombinatorics.Det.det_eq_cols, AlgebraicCombinatorics.Det.det_scale_col, AlgebraicCombinatorics.Det.det_add_col, AlgebraicCombinatorics.Det.det_add_smul_col, AlgebraicCombinatorics.Det.det_col_add}
\leantarget
\leanok
Theorem~\ref{thm.det.rowop} also holds if we replace ``row'' by ``column''
throughout it.
\end{theorem}

\begin{proof}
\leanok
Theorem~\ref{thm.det.transp} shows that the determinant of a matrix does not
change when we replace it by its transpose; however, the rows of this
transpose $A^T$ are the transposes of the columns of $A$. Thus,
Theorem~\ref{thm.det.colop} follows by applying Theorem~\ref{thm.det.rowop}
to the transposes of all the matrices involved.
\end{proof}

\begin{corollary}
\label{cor.det.sig-row-col}
\lean{AlgebraicCombinatorics.Det.det_permute_rows, AlgebraicCombinatorics.Det.det_permute_cols}
\leantarget
\leanok
Let $n \in \mathbb{N}$. Let $A \in K^{n \times n}$ and $\tau \in S_n$. Then,
\begin{equation}
  \det\left( \left( A_{\tau(i), j} \right)_{1 \leq i \leq n,\; 1 \leq j \leq n} \right)
  = (-1)^{\tau} \cdot \det A
  \label{eq.cor.det.sig-row-col.row}
\end{equation}
and
\begin{equation}
  \det\left( \left( A_{i, \tau(j)} \right)_{1 \leq i \leq n,\; 1 \leq j \leq n} \right)
  = (-1)^{\tau} \cdot \det A.
  \label{eq.cor.det.sig-row-col.col}
\end{equation}
\end{corollary}

\begin{proof}
\leanok
Let us first prove \eqref{eq.cor.det.sig-row-col.col}.

The definition of $\det A$ yields
\begin{align*}
  \det A
  &= \sum_{\sigma \in S_n} (-1)^{\sigma} \prod_{i=1}^{n} A_{i, \sigma(i)}.
\end{align*}
The definition of $\det\left( \left( A_{i, \tau(j)} \right)_{1 \leq i \leq n,\;
1 \leq j \leq n} \right)$ yields
\begin{align*}
  \det\left( \left( A_{i, \tau(j)} \right)_{1 \leq i \leq n,\; 1 \leq j \leq n} \right)
  &= \sum_{\sigma \in S_n} (-1)^{\sigma} \prod_{i=1}^{n} A_{i, \tau(\sigma(i))}.
\end{align*}
However, $S_n$ is a group. Thus, the map $\sigma \mapsto \tau^{-1} \sigma$
is a bijection on $S_n$. Substituting $\tau^{-1} \sigma$ for $\sigma$, we obtain
(using Proposition~\ref{prop.perm.sign.props}~\textbf{(d)} and \textbf{(f)})
\begin{align*}
  &\sum_{\sigma \in S_n} (-1)^{\sigma} \prod_{i=1}^{n} A_{i, \tau(\sigma(i))} \\
  &= \sum_{\sigma \in S_n}
     \underbrace{(-1)^{\tau^{-1} \sigma}}_{
       = (-1)^{\tau^{-1}} \cdot (-1)^{\sigma}}
     \prod_{i=1}^{n}
     \underbrace{A_{i, \tau((\tau^{-1} \sigma)(i))}}_{= A_{i, \sigma(i)}} \\
  &= \underbrace{(-1)^{\tau^{-1}}}_{= (-1)^{\tau}} \cdot
     \underbrace{\sum_{\sigma \in S_n} (-1)^{\sigma} \prod_{i=1}^{n} A_{i, \sigma(i)}}_{= \det A} \\
  &= (-1)^{\tau} \cdot \det A.
\end{align*}
This proves \eqref{eq.cor.det.sig-row-col.col}.

Now, \eqref{eq.cor.det.sig-row-col.row} follows by applying
\eqref{eq.cor.det.sig-row-col.col} to $A^T$ instead of $A$ and using
Theorem~\ref{thm.det.transp}.
\end{proof}

\begin{theorem}[Multiplicativity of the determinant]
\label{thm.det.detAB}
\lean{AlgebraicCombinatorics.Det.det_mul'}
\leantarget
\leanok
Let $n \in \mathbb{N}$. Let $A, B \in K^{n \times n}$ be two $n \times n$-matrices.
Then,
\[
  \det(AB) = \det A \cdot \det B.
\]
\end{theorem}

\begin{proof}
\leanok
See \cite[Theorem B.17]{Strick13} or \cite[Theorem 6.23]{detnotes} or
\cite[\S 5]{Zeilbe} or \cite[Lemma 4.7.5 (1)]{Ford21} or
\cite[Theorem 5.7]{Laue-det}.
\end{proof}

\begin{corollary}
\label{cor.det.scale-row-col}
\lean{AlgebraicCombinatorics.Det.det_scale_rows, AlgebraicCombinatorics.Det.det_scale_cols}
\leantarget
\leanok
Let $n \in \mathbb{N}$. Let $A \in K^{n \times n}$ and
$d_1, d_2, \ldots, d_n \in K$. Then,
\begin{equation}
  \det\left( \left( d_i A_{i,j} \right)_{1 \leq i \leq n,\; 1 \leq j \leq n} \right)
  = d_1 d_2 \cdots d_n \cdot \det A
  \label{eq.cor.det.scale-row-col.row}
\end{equation}
and
\begin{equation}
  \det\left( \left( d_j A_{i,j} \right)_{1 \leq i \leq n,\; 1 \leq j \leq n} \right)
  = d_1 d_2 \cdots d_n \cdot \det A.
  \label{eq.cor.det.scale-row-col.col}
\end{equation}
\end{corollary}

\begin{proof}
\leanok
The definition of a determinant yields
\[
  \det A = \sum_{\sigma \in S_n} (-1)^{\sigma}
  A_{1,\sigma(1)} A_{2,\sigma(2)} \cdots A_{n,\sigma(n)}
\]
and
\begin{align*}
  &\det\left( \left( d_i A_{i,j} \right)_{1 \leq i \leq n,\; 1 \leq j \leq n} \right) \\
  &= \sum_{\sigma \in S_n} (-1)^{\sigma}
     \underbrace{(d_1 A_{1,\sigma(1)})(d_2 A_{2,\sigma(2)}) \cdots (d_n A_{n,\sigma(n)})}_{
       = (d_1 d_2 \cdots d_n)(A_{1,\sigma(1)} A_{2,\sigma(2)} \cdots A_{n,\sigma(n)})} \\
  &= d_1 d_2 \cdots d_n \cdot
     \underbrace{\sum_{\sigma \in S_n} (-1)^{\sigma}
       A_{1,\sigma(1)} A_{2,\sigma(2)} \cdots A_{n,\sigma(n)}}_{= \det A} \\
  &= d_1 d_2 \cdots d_n \cdot \det A
\end{align*}
and
\begin{align*}
  &\det\left( \left( d_j A_{i,j} \right)_{1 \leq i \leq n,\; 1 \leq j \leq n} \right) \\
  &= \sum_{\sigma \in S_n} (-1)^{\sigma}
     \underbrace{(d_{\sigma(1)} A_{1,\sigma(1)})(d_{\sigma(2)} A_{2,\sigma(2)})
       \cdots (d_{\sigma(n)} A_{n,\sigma(n)})}_{
       = (d_{\sigma(1)} d_{\sigma(2)} \cdots d_{\sigma(n)})
         (A_{1,\sigma(1)} A_{2,\sigma(2)} \cdots A_{n,\sigma(n)})} \\
  &= \sum_{\sigma \in S_n} (-1)^{\sigma}
     \underbrace{(d_{\sigma(1)} d_{\sigma(2)} \cdots d_{\sigma(n)})}_{\substack{
       = d_1 d_2 \cdots d_n \\
       \text{(since } \sigma \text{ is a permutation of } [n] \text{)}}}
     (A_{1,\sigma(1)} A_{2,\sigma(2)} \cdots A_{n,\sigma(n)}) \\
  &= d_1 d_2 \cdots d_n \cdot \det A.
\end{align*}
\end{proof}

\subsection{Permutation images and submatrix determinants}

\begin{definition}
\label{def.det.imageFinset}
\lean{PermFinset.imageFinset}
\leanhelper
\leanok
The \emph{image} of a finite subset $P \subseteq [n]$ under a permutation
$\sigma \in S_n$ is $\sigma(P) = \{\sigma(i) \mid i \in P\}$.
\end{definition}

\begin{lemma}
\label{lem.det.imageFinset-card}
\lean{PermFinset.imageFinset_card}
\leanhelper
\leanok
The image of a finset under a permutation has the same cardinality:
$|\sigma(P)| = |P|$.
\end{lemma}

\begin{proof}
\leanok
Since $\sigma$ is injective, the image has the same cardinality as the original set.
\end{proof}

\begin{lemma}
\label{lem.det.imageFinset-compl}
\lean{PermFinset.imageFinset_compl}
\leanhelper
\leanok
The image of a complement equals the complement of the image:
$\sigma(P^c) = \sigma(P)^c$.
\end{lemma}

\begin{proof}
\leanok
By the injectivity of $\sigma$: $x \in \sigma(P^c)$ iff $\sigma^{-1}(x) \notin P$
iff $x \notin \sigma(P)$.
\end{proof}

\begin{lemma}
\label{lem.det.imageFinset-inv}
\lean{PermFinset.imageFinset_inv}
\leanhelper
\leanok
If $\sigma(P) = Q$, then $\sigma^{-1}(Q) = P$.
\end{lemma}

\begin{proof}
\leanok
Apply $\sigma^{-1}$ to both sides of $\sigma(P) = Q$ and use
$\sigma^{-1} \circ \sigma = \mathrm{id}$.
\end{proof}

\begin{definition}
\label{def.det.permsMapping}
\lean{PermFinset.permsMapping}
\leanhelper
\leanok
The set of permutations mapping $P$ to $Q$:
$\mathrm{permsMapping}(P, Q) = \{\sigma \in S_n \mid \sigma(P) = Q\}$.
\end{definition}

\begin{lemma}
\label{lem.det.permsMapping-empty}
\lean{PermFinset.permsMapping_empty_of_card_ne}
\leanhelper
\leanok
If $|P| \neq |Q|$, then no permutation maps $P$ to $Q$:
$\mathrm{permsMapping}(P, Q) = \emptyset$.
\end{lemma}

\begin{proof}
\leanok
If $\sigma(P) = Q$, then $|Q| = |\sigma(P)| = |P|$
by Lemma~\ref{lem.det.imageFinset-card}, contradicting $|P| \neq |Q|$.
\end{proof}

\begin{definition}
\label{def.det.submatrixDet}
\lean{PermFinset.submatrixDet}
\leanhelper
\leanok
The \emph{submatrix determinant} of an $n \times n$ matrix $A$ with row set $P$
and column set $Q$ (both subsets of $[n]$) is
\[
\mathrm{submatrixDet}(A, P, Q) =
\begin{cases}
\det(A[P, Q]) & \text{if } |P| = |Q|, \\
0 & \text{otherwise,}
\end{cases}
\]
where $A[P, Q]$ is the submatrix of $A$ with rows indexed by $P$ and columns
indexed by $Q$, in the order determined by their natural ordering.
\end{definition}

\begin{lemma}
\label{lem.det.submatrixDet-card-eq}
\lean{PermFinset.submatrixDet_of_card_eq}
\leanhelper
\leanok
When $|P| = |Q|$, the submatrix determinant equals the actual determinant
of the submatrix $A[P, Q]$.
\end{lemma}

\begin{proof}
\leanok
Immediate from the definitions, using the hypothesis $|P| = |Q|$.
\end{proof}

\begin{lemma}
\label{lem.det.submatrixDet-card-ne}
\lean{PermFinset.submatrixDet_of_card_ne}
\leanhelper
\leanok
When $|P| \neq |Q|$, the submatrix determinant is $0$.
\end{lemma}

\begin{proof}
\leanok
Immediate from the definitions, using the hypothesis $|P| \neq |Q|$.
\end{proof}

\begin{lemma}
\label{lem.det.submatrixDet-empty}
\lean{PermFinset.submatrixDet_empty}
\leanhelper
\leanok
The submatrix determinant of the empty row and column sets is $1$:
$\mathrm{submatrixDet}(A, \emptyset, \emptyset) = 1$.
\end{lemma}

\begin{proof}
\leanok
Both sets are empty (so $|P| = |Q| = 0$), and the determinant of the
$0 \times 0$ submatrix is $1$.
\end{proof}
